% ============================================================
\chapter{Reproductibilit\'e en Recherche --- Standards et Bonnes Pratiques}
\label{ch:reproductibilite}
\index{reproductibilit\'e}
% ============================================================

\section{La crise de la reproductibilit\'e en ML}
\label{sec:crise-repro}
\index{reproductibilit\'e!crise}

La \textbf{crise de la reproductibilit\'e}\index{crise de la reproductibilit\'e}
est un probl\`eme majeur en recherche scientifique, et le machine learning
n'y \'echappe pas. Plusieurs \'etudes ont montr\'e que la majorit\'e des
r\'esultats publi\'es en ML sont difficiles, voire impossibles, \`a reproduire.

\begin{attention}[title=\textbf{Chiffres alarmants}]
\begin{itemize}
  \item Selon une enqu\^ete de \emph{Nature} (2016), plus de 70\,\% des
    chercheurs ont \'echou\'e \`a reproduire les r\'esultats d'un autre
    chercheur.
  \item En NLP, Dodge et al.\ (2019) ont montr\'e que le simple choix de la
    graine al\'eatoire pouvait faire varier la performance de 1 \`a 5 points
    de F1-score.
  \item En vision par ordinateur, de nombreux r\'esultats <<~\'etat de l'art~>>
    ne sont pas reproductibles sans acc\`es au code et aux donn\'ees originaux.
\end{itemize}
\end{attention}

Les causes sp\'ecifiques au ML sont multiples~:

\begin{itemize}
  \item \textbf{Al\'eatoire non contr\^ol\'e}~: initialisation des poids,
    ordre des batchs, augmentation de donn\'ees.
  \item \textbf{D\'ependances logicielles}~: versions de CUDA, cuDNN, PyTorch
    qui influencent les r\'esultats num\'eriques.
  \item \textbf{Donn\'ees non partag\'ees}~: datasets propri\'etaires,
    pr\'etraitements non document\'es.
  \item \textbf{Hyperparamm\`etres cach\'es}~: param\`etres non mentionn\'es
    dans l'article (learning rate scheduler, warmup, etc.).
  \item \textbf{Co\^ut computationnel}~: entra\^inements n\'ecessitant des
    centaines de GPU-heures, impossibles \`a re-ex\'ecuter.
\end{itemize}

\section{Niveaux de reproductibilit\'e}
\label{sec:niveaux-repro}
\index{reproductibilit\'e!niveaux}

Il est essentiel de distinguer trois niveaux de reproductibilit\'e, souvent
confondus~:

\begin{definition}[R\'ep\'etabilit\'e]
La \textbf{r\'ep\'etabilit\'e}\index{r\'ep\'etabilit\'e} (\emph{repeatability})
est la capacit\'e d'obtenir les m\^emes r\'esultats en r\'e-ex\'ecutant le
m\^eme code, sur les m\^emes donn\'ees, dans le m\^eme environnement.
C'est le niveau le plus basique.
\end{definition}

\begin{definition}[Reproductibilit\'e]
La \textbf{reproductibilit\'e}\index{reproductibilit\'e!d\'efinition}
(\emph{reproducibility}) est la capacit\'e pour une \'equipe ind\'ependante
d'obtenir les m\^emes r\'esultats en utilisant le code et les donn\'ees
fournis par les auteurs, mais dans leur propre environnement.
\end{definition}

\begin{definition}[R\'eplicabilit\'e]
La \textbf{r\'eplicabilit\'e}\index{r\'eplicabilit\'e} (\emph{replicability})
est la capacit\'e d'obtenir des r\'esultats similaires avec une
impl\'ementation ind\'ependante, sur des donn\'ees diff\'erentes. C'est le
niveau le plus exigeant, qui valide la g\'en\'eralisabilit\'e de la m\'ethode.
\end{definition}

\begin{center}
\begin{tikzpicture}[
  level/.style={rectangle, draw, rounded corners, minimum width=11cm,
    minimum height=1.6cm, align=center, font=\small},
  arr/.style={-{Stealth[length=2.5mm]}, thick, steelblue}
]
  \node[level, fill=green!10] (rep) {
    \textbf{R\'ep\'etabilit\'e (Repeatability)}\\
    M\^eme \'equipe, m\^eme code, m\^emes donn\'ees, m\^eme environnement
  };
  \node[level, fill=yellow!10, above=0.6cm of rep] (repro) {
    \textbf{Reproductibilit\'e (Reproducibility)}\\
    \'Equipe diff\'erente, m\^eme code, m\^emes donn\'ees, environnement diff\'erent
  };
  \node[level, fill=orange!10, above=0.6cm of repro] (repli) {
    \textbf{R\'eplicabilit\'e (Replicability)}\\
    \'Equipe diff\'erente, code diff\'erent, donn\'ees diff\'erentes
  };

  \draw[arr] (rep) -- (repro);
  \draw[arr] (repro) -- (repli);

  \node[right=0.3cm of rep, font=\scriptsize, text=green!50!black] {Facile};
  \node[right=0.3cm of repli, font=\scriptsize, text=red!60!black] {Difficile};
\end{tikzpicture}
\end{center}

\begin{remarque}
En pratique, m\^eme la r\'ep\'etabilit\'e n'est pas garantie en ML~: les
op\'erations GPU non d\'eterministes (atomicAdd, cuDNN autotuner) peuvent
produire des r\'esultats l\'eg\`erement diff\'erents d'une ex\'ecution \`a
l'autre sur le m\^eme mat\'eriel.
\end{remarque}

\section{Gestion des graines al\'eatoires}
\label{sec:random-seeds}
\index{graine al\'eatoire}
\index{reproductibilit\'e!graines al\'eatoires}

Le contr\^ole de l'al\'eatoire est la premi\`ere \'etape vers la
r\'ep\'etabilit\'e. En Python, plusieurs sources d'al\'eatoire coexistent
et doivent toutes \^etre fix\'ees.

\begin{codeblock}[title=\textbf{Configuration compl\`ete de reproductibilit\'e}]
\begin{minted}{python}
"""Reproducibility utilities for PyTorch projects."""
import os
import random
import numpy as np
import torch


def set_seed(seed: int = 42, deterministic: bool = True) -> None:
    """Set all random seeds for reproducibility.

    Args:
        seed: Random seed value.
        deterministic: If True, enforce deterministic algorithms
            (may reduce performance).
    """
    # Python built-in random
    random.seed(seed)

    # Numpy
    np.random.seed(seed)

    # PyTorch CPU
    torch.manual_seed(seed)

    # PyTorch GPU (all devices)
    torch.cuda.manual_seed(seed)
    torch.cuda.manual_seed_all(seed)

    # Hash seed for Python dicts, sets, etc.
    os.environ["PYTHONHASHSEED"] = str(seed)

    if deterministic:
        # Force deterministic algorithms
        torch.use_deterministic_algorithms(True)

        # Disable cuDNN benchmark (non-deterministic)
        torch.backends.cudnn.benchmark = False
        torch.backends.cudnn.deterministic = True

        # Handle CUBLAS non-determinism
        os.environ["CUBLAS_WORKSPACE_CONFIG"] = ":4096:8"


def seed_worker(worker_id: int) -> None:
    """Seed function for DataLoader workers.

    Usage:
        DataLoader(..., worker_init_fn=seed_worker)
    """
    worker_seed = torch.initial_seed() % 2**32
    np.random.seed(worker_seed)
    random.seed(worker_seed)


# Usage
set_seed(42)

# DataLoader with reproducible workers
g = torch.Generator()
g.manual_seed(42)

train_loader = torch.utils.data.DataLoader(
    dataset,
    batch_size=32,
    shuffle=True,
    num_workers=4,
    worker_init_fn=seed_worker,
    generator=g,
)
\end{minted}
\end{codeblock}

\begin{attention}[title=\textbf{D\'eterminisme CUDA~: le prix \`a payer}]
Activer \cli{torch.use\_deterministic\_algorithms(True)} peut~:
\begin{itemize}
  \item Ralentir l'entra\^inement de 10 \`a 20\,\%.
  \item Lever des erreurs pour certaines op\'erations qui n'ont pas
    d'impl\'ementation d\'eterministe (ex.~: \cli{scatter\_add},
    certaines convolutions).
  \item Ne pas \^etre suffisant entre versions de CUDA ou de mat\'eriel
    diff\'erents.
\end{itemize}
\textbf{Recommandation}~: utiliser le d\'eterminisme pour le d\'ebogage et
la validation, mais accepter une l\'eg\`ere variabilit\'e en production.
\end{attention}

\begin{formules}[title=\textbf{Sources d'al\'eatoire \`a contr\^oler}]
\begin{tabular}{@{}lll@{}}
\toprule
\textbf{Source} & \textbf{Commande} & \textbf{Impact} \\
\midrule
Python \cli{random} & \cli{random.seed(s)} & Shuffling, sampling \\
NumPy & \cli{np.random.seed(s)} & Initialisation, bruit \\
PyTorch CPU & \cli{torch.manual\_seed(s)} & Poids initiaux \\
PyTorch GPU & \cli{torch.cuda.manual\_seed\_all(s)} & Op\'erations CUDA \\
cuDNN & \cli{cudnn.deterministic = True} & Convolutions \\
CUBLAS & \cli{CUBLAS\_WORKSPACE\_CONFIG} & Produits matriciels \\
DataLoader & \cli{worker\_init\_fn} & Workers parall\`eles \\
PYTHONHASHSEED & \cli{os.environ[...] = str(s)} & Ordre des dicts \\
\bottomrule
\end{tabular}
\end{formules}

\section{Documentation des exp\'eriences}
\label{sec:doc-experiments}
\index{model card}
\index{datasheet}

Au-del\`a du code, la documentation est essentielle pour la
reproductibilit\'e. Deux standards \'emergent dans la communaut\'e ML.

\subsection{Model Cards}
\index{model card}

Les \textbf{Model Cards} (Mitchell et al., 2019) sont des documents
structur\'es qui accompagnent un mod\`ele publi\'e.

\begin{bonnepratique}[title=\textbf{Contenu d'une Model Card}]
\begin{enumerate}
  \item \textbf{D\'etails du mod\`ele}~: architecture, version, auteurs,
    licence.
  \item \textbf{Usage pr\'evu}~: cas d'utilisation primaires et hors-scope.
  \item \textbf{M\'etriques}~: m\'etriques choisies et justification.
  \item \textbf{Donn\'ees d'entra\^inement}~: description, taille, source.
  \item \textbf{Donn\'ees d'\'evaluation}~: dataset de test, m\'ethodologie.
  \item \textbf{Analyses \'ethiques}~: biais, \'equit\'e, limites connues.
  \item \textbf{R\'esultats quantitatifs}~: par sous-groupe
    d\'emographique si applicable.
  \item \textbf{Consid\'erations environnementales}~: empreinte carbone,
    co\^ut computationnel.
\end{enumerate}
\end{bonnepratique}

\begin{codeblock}[title=\textbf{Model Card en YAML (extrait)}]
\begin{minted}{yaml}
model_card:
  model_details:
    name: "sentiment-classifier-v2"
    version: "2.1.0"
    architecture: "BERT-base fine-tuned"
    parameters: 110_000_000
    license: "Apache-2.0"
    training_date: "2025-01-15"
    training_cost: "8 GPU-hours (A100)"

  intended_use:
    primary: "Sentiment analysis of product reviews"
    out_of_scope: "Medical text, legal documents"

  metrics:
    - name: "F1-score (macro)"
      value: 0.892
      confidence_interval: [0.885, 0.899]
    - name: "Accuracy"
      value: 0.904

  training_data:
    name: "Amazon Reviews (2023)"
    size: "3.6M reviews"
    preprocessing: "Lowercased, truncated to 512 tokens"

  ethical_considerations:
    bias: "Lower performance on non-English reviews"
    fairness_evaluation: "Tested across 5 product categories"
\end{minted}
\end{codeblock}

\subsection{Datasheets for Datasets}
\index{datasheets for datasets}

Les \textbf{Datasheets for Datasets} (Gebru et al., 2021) documentent les
datasets de mani\`ere standardis\'ee~:

\begin{itemize}
  \item \textbf{Motivation}~: pourquoi le dataset a-t-il \'et\'e cr\'e\'e~?
  \item \textbf{Composition}~: quelles instances contient-il~? Quelle est
    la distribution~?
  \item \textbf{Collecte}~: comment les donn\'ees ont-elles \'et\'e
    collect\'ees~? Par qui~?
  \item \textbf{Pr\'etraitement}~: quelles transformations ont \'et\'e
    appliqu\'ees~?
  \item \textbf{Distribution}~: comment le dataset est-il distribu\'e~?
    Sous quelle licence~?
  \item \textbf{Maintenance}~: qui maintient le dataset~? Comment signaler
    des probl\`emes~?
\end{itemize}

\section{Archivage et partage des artefacts}
\label{sec:archivage}
\index{archivage}
\index{Zenodo}
\index{Papers with Code}

La reproductibilit\'e exige que le code, les donn\'ees et les mod\`eles
soient archiv\'es de mani\`ere p\'erenne et accessible.

\begin{formules}[title=\textbf{Plateformes d'archivage}]
\begin{tabular}{@{}llp{6cm}@{}}
\toprule
\textbf{Plateforme} & \textbf{Type} & \textbf{Cas d'usage} \\
\midrule
Zenodo & Code + donn\'ees & Archive p\'erenne avec DOI, int\'egration GitHub,
  jusqu'\`a 50\,Go par d\'ep\^ot \\
Papers with Code & Code + r\'esultats & Lier articles, code et benchmarks~;
  leaderboards communautaires \\
Hugging Face Hub & Mod\`eles + datasets & H\'ebergement de mod\`eles et
  datasets avec versionnement \\
GitHub Releases & Code & Snapshots versionn\'es du code \\
DVC (remote) & Donn\'ees + mod\`eles & Versionnement de fichiers volumineux
  sur S3, GCS, Azure \\
\bottomrule
\end{tabular}
\end{formules}

\begin{codeblock}[title=\textbf{Archivage automatique avec Zenodo + GitHub}]
\begin{minted}{bash}
# 1. Activer l'integration Zenodo-GitHub
#    -> https://zenodo.org/account/settings/github/

# 2. Creer une release GitHub
git tag -a v1.0.0 -m "Camera-ready version for NeurIPS 2025"
git push origin v1.0.0

# 3. Zenodo cree automatiquement un DOI pour la release

# 4. Ajouter le badge DOI dans le README
# [![DOI](https://zenodo.org/badge/DOI/10.5281/zenodo.XXXXX.svg)]
\end{minted}
\end{codeblock}

\begin{codeblock}[title=\textbf{Publication sur Hugging Face Hub}]
\begin{minted}{python}
from huggingface_hub import HfApi, ModelCard

api = HfApi()

# Upload model
api.upload_folder(
    folder_path="./model_checkpoint",
    repo_id="username/sentiment-classifier-v2",
    repo_type="model",
)

# Upload dataset
api.upload_folder(
    folder_path="./processed_data",
    repo_id="username/amazon-reviews-2023",
    repo_type="dataset",
)

# Create and push model card
card = ModelCard.load("username/sentiment-classifier-v2")
card.data.language = "en"
card.data.license = "apache-2.0"
card.data.datasets = ["username/amazon-reviews-2023"]
card.push_to_hub("username/sentiment-classifier-v2")
\end{minted}
\end{codeblock}

\section{Checklist de reproductibilit\'e}
\label{sec:checklist}
\index{NeurIPS!checklist}
\index{reproductibilit\'e!checklist}

La conf\'erence NeurIPS exige depuis 2019 une checklist de
reproductibilit\'e. Voici une version adapt\'ee pour tout projet ML~:

\begin{bonnepratique}[title=\textbf{Checklist de reproductibilit\'e ML}]
\textbf{Code et environnement~:}
\begin{itemize}
  \item[$\square$] Code source disponible publiquement (ou en suppl\'ement)
  \item[$\square$] Fichier de d\'ependances complet avec versions \'epingl\'ees
  \item[$\square$] Instructions d'installation test\'ees sur une machine vierge
  \item[$\square$] Script unique pour reproduire tous les r\'esultats
  \item[$\square$] Version de Python, CUDA, cuDNN document\'ee
\end{itemize}
\textbf{Donn\'ees~:}
\begin{itemize}
  \item[$\square$] Donn\'ees disponibles ou instructions pour les obtenir
  \item[$\square$] Pr\'etraitement enti\`erement script\'e (pas d'\'etapes manuelles)
  \item[$\square$] Splits train/val/test fix\'es et partag\'es
  \item[$\square$] Datasheet ou documentation du dataset
\end{itemize}
\textbf{Exp\'eriences~:}
\begin{itemize}
  \item[$\square$] Tous les hyperparamm\`etres rapport\'es
  \item[$\square$] Graines al\'eatoires fix\'ees et document\'ees
  \item[$\square$] R\'esultats moyenn\'es sur plusieurs runs ($\geq 3$)
  \item[$\square$] Barres d'erreur ou intervalles de confiance rapport\'es
  \item[$\square$] Temps d'entra\^inement et ressources document\'es
\end{itemize}
\textbf{Mod\`ele~:}
\begin{itemize}
  \item[$\square$] Architecture compl\`etement d\'ecrite
  \item[$\square$] Nombre de param\`etres rapport\'e
  \item[$\square$] Model card incluse
  \item[$\square$] Poids du mod\`ele disponibles (si applicable)
\end{itemize}
\end{bonnepratique}

\section{Comparaison des outils pour la recherche reproductible}
\label{sec:tools-repro}

\begin{center}
\begin{tabular}{@{}lp{3cm}p{3cm}p{3.5cm}@{}}
\toprule
\textbf{Outil} & \textbf{Points forts} & \textbf{Limites} &
  \textbf{Cas d'usage} \\
\midrule
MLflow & Suivi d'exp\'eriences, registre de mod\`eles & Interface parfois
  lourde & Projets d'\'equipe, production \\
DVC & Versionnement donn\'ees, pipelines & Courbe d'apprentissage &
  Gros datasets, pipelines complexes \\
W\&B & Visualisation, collaboration & SaaS (donn\'ees sur le cloud) &
  Recherche acad\'emique, prototypage \\
Hydra & Configuration flexible, multi-run & Code sp\'ecifique &
  Balayage d'hyperparamm\`etres \\
Sacred & L\'eger, MongoDB & Peu maintenu & Petits projets acad\'emiques \\
Guild AI & CLI, comparaison de runs & Communaut\'e r\'eduite &
  Exp\'erimentation locale \\
\bottomrule
\end{tabular}
\end{center}

\section{Bonnes pratiques compl\'ementaires}
\label{sec:bonnes-pratiques-repro}

\begin{bonnepratique}[title=\textbf{R\`egles d'or de la reproductibilit\'e}]
\begin{enumerate}
  \item \textbf{Versionner tout}~: code (Git), donn\'ees (DVC), environnement
    (\file{requirements.txt}), configuration (YAML).
  \item \textbf{Automatiser tout}~: un seul \cli{make reproduce} doit
    g\'en\'erer tous les r\'esultats.
  \item \textbf{Documenter les choix}~: pourquoi ce learning rate~? Pourquoi
    cette architecture~? Notez les d\'ecisions dans un journal d'exp\'eriences.
  \item \textbf{Rapporter la variabilit\'e}~: toujours moyenne $\pm$
    \'ecart-type sur $n \geq 3$ runs.
  \item \textbf{Conteneuriser}~: un \file{Dockerfile} garantit un
    environnement identique.
  \item \textbf{Archiver}~: d\'eposez code et donn\'ees sur Zenodo pour
    obtenir un DOI p\'erenne.
\end{enumerate}
\end{bonnepratique}

\begin{antipattern}[title=\textbf{Anti-patterns de la reproductibilit\'e}]
\begin{itemize}
  \item <<~Les r\'esultats sont dans le notebook~>> sans script de
    reproduction.
  \item Rapporter le meilleur run sur 100 sans mentionner les 99 autres.
  \item Utiliser \cli{random\_state=42} partout sans justifier ni
    tester d'autres valeurs.
  \item Ne pas versionner les donn\'ees de pr\'etraitement.
  \item <<~Code disponible sur demande~>> (spoiler~: il ne l'est jamais).
  \item Modifier le code apr\`es soumission sans mettre \`a jour le d\'ep\^ot.
\end{itemize}
\end{antipattern}

\section{Mini-projet~: rendre le projet de cours reproductible}
\label{sec:mini-project-ch11}

\begin{miniprojet}[title=\textbf{Mini-projet 11~: Reproductibilit\'e compl\`ete}]
Prenez le projet complet d\'evelopp\'e au fil des chapitres pr\'ec\'edents et
rendez-le enti\`erement reproductible~:
\begin{enumerate}
  \item \textbf{Graines al\'eatoires}~: impl\'ementez la fonction
    \cli{set\_seed()} compl\`ete et appelez-la au d\'ebut de chaque script.
  \item \textbf{Conteneurisation}~: cr\'eez un \file{Dockerfile} qui
    installe toutes les d\'ependances et reproduit l'entra\^inement.
  \item \textbf{Pipeline DVC}~: d\'efinissez un \file{dvc.yaml} avec les
    \'etapes preprocess $\to$ train $\to$ evaluate.
  \item \textbf{Documentation}~:
    \begin{itemize}
      \item R\'edigez une Model Card pour votre mod\`ele final.
      \item R\'edigez une Datasheet pour votre dataset.
    \end{itemize}
  \item \textbf{Script de reproduction}~: cr\'eez un \file{reproduce.sh}
    qui clone le d\'ep\^ot, construit le conteneur, et ex\'ecute le pipeline
    complet.
  \item \textbf{Validation}~: demandez \`a un coll\`egue de reproduire vos
    r\'esultats en suivant uniquement le README. Notez les probl\`emes
    rencontr\'es et corrigez-les.
  \item \textbf{Archivage}~: cr\'eez une release GitHub et archivez-la
    sur Zenodo.
\end{enumerate}
\end{miniprojet}

\begin{codeblock}[title=\textbf{reproduce.sh --- Script de reproduction compl\`ete}]
\begin{minted}{bash}
#!/bin/bash
set -euo pipefail

echo "=== Reproduction complete du projet ==="

# 1. Verifier les prerequis
command -v docker >/dev/null 2>&1 || {
    echo "Docker requis. Installation: https://docs.docker.com/install/"
    exit 1
}

# 2. Construire l'image Docker
echo "[1/4] Construction de l'image Docker..."
docker build -t ml-project:reproduce .

# 3. Telecharger les donnees versionnees
echo "[2/4] Telechargement des donnees (DVC)..."
docker run --rm -v "$(pwd)/data:/app/data" ml-project:reproduce \
    dvc pull

# 4. Executer le pipeline complet
echo "[3/4] Execution du pipeline..."
docker run --rm \
    -v "$(pwd)/data:/app/data" \
    -v "$(pwd)/results:/app/results" \
    ml-project:reproduce \
    dvc repro

# 5. Verifier les resultats
echo "[4/4] Verification des metriques..."
docker run --rm \
    -v "$(pwd)/results:/app/results" \
    ml-project:reproduce \
    python src/verify_results.py \
        --expected-accuracy 0.892 \
        --tolerance 0.005

echo "=== Reproduction terminee avec succes ==="
\end{minted}
\end{codeblock}

\section{Exercices}
\label{sec:exercises-ch11}

\begin{exercice}[$\bigstar$ --- Audit de reproductibilit\'e]
Choisissez un article r\'ecent sur \emph{Papers with Code}. V\'erifiez~:
\begin{enumerate}
  \item Le code est-il disponible~?
  \item Les d\'ependances sont-elles document\'ees~?
  \item Les hyperparamm\`etres sont-ils tous rapport\'es~?
  \item Les graines al\'eatoires sont-elles fix\'ees~?
  \item Pouvez-vous ex\'ecuter le code avec succ\`es~?
\end{enumerate}
R\'edigez un rapport d'audit de 2 pages en utilisant la checklist de la
section~\ref{sec:checklist}.
\end{exercice}

\begin{exercice}[$\bigstar\bigstar$ --- Impact de l'al\'eatoire]
Entra\^inez un r\'eseau de neurones (par exemple un CNN sur CIFAR-10) avec
10 graines al\'eatoires diff\'erentes. Pour chaque configuration~:
\begin{enumerate}
  \item Mesurez l'accuracy finale sur le test set.
  \item Calculez la moyenne, l'\'ecart-type, le min et le max.
  \item Comparez les r\'esultats avec et sans
    \cli{torch.use\_deterministic\_algorithms(True)}.
  \item Tracez un box plot des performances.
  \item Mesurez l'impact sur le temps d'entra\^inement.
\end{enumerate}
\end{exercice}

\begin{exercice}[$\bigstar\bigstar\bigstar$ --- Pipeline reproductible de bout en bout]
Construisez un projet ML enti\`erement reproductible~:
\begin{enumerate}
  \item Choisissez un probl\`eme de classification (texte ou image).
  \item Impl\'ementez le code avec la structure modulaire du
    chapitre~\ref{ch:introduction}.
  \item Fixez toutes les sources d'al\'eatoire.
  \item Cr\'eez un pipeline DVC complet.
  \item Conteneurisez avec Docker.
  \item R\'edigez une Model Card et une Datasheet.
  \item \'Ecrivez un \file{reproduce.sh} qui automatise tout.
  \item Faites reproduire le projet par un coll\`egue et it\'erez
    sur les probl\`emes rencontr\'es.
  \item Archivez le tout sur Zenodo et obtenez un DOI.
  \item Appliquez la checklist NeurIPS et v\'erifiez que toutes les
    cases sont coch\'ees.
\end{enumerate}
\end{exercice}

\section{Aide-m\'emoire}

\begin{cheatsheet}[title=\textbf{R\'esum\'e du chapitre 11}]
\begin{tabular}{@{}p{4cm}p{8.5cm}@{}}
\toprule
\textbf{Concept} & \textbf{Point cl\'e} \\
\midrule
R\'ep\'etabilit\'e & M\^eme code, m\^emes donn\'ees, m\^eme environnement
  $\to$ m\^emes r\'esultats \\
Reproductibilit\'e & M\^eme code et donn\'ees, environnement
  diff\'erent $\to$ m\^emes r\'esultats \\
R\'eplicabilit\'e & Code et donn\'ees diff\'erents $\to$ conclusions
  similaires \\
Graines al\'eatoires & Fixer Python, NumPy, PyTorch, CUDA, DataLoader \\
D\'eterminisme GPU & \cli{torch.use\_deterministic\_algorithms(True)},
  ralentissement de 10--20\,\% \\
Model Card & Documentation structur\'ee du mod\`ele (architecture, biais,
  limites) \\
Datasheet & Documentation structur\'ee du dataset (collecte, composition,
  licence) \\
Archivage & Zenodo (DOI), Hugging Face Hub, Papers with Code \\
Checklist & Code, donn\'ees, hyperparamm\`etres, variabilit\'e, ressources \\
Automatisation & \file{reproduce.sh} + Docker + DVC pour reproduction
  compl\`ete \\
\bottomrule
\end{tabular}
\end{cheatsheet}
